% =====================================================================
%  Capitolo 1 — Introduzione
%  Incluso da main.tex con % =====================================================================
%  Capitolo 1 — Introduzione
%  Incluso da main.tex con % =====================================================================
%  Capitolo 1 — Introduzione
%  Incluso da main.tex con % =====================================================================
%  Capitolo 1 — Introduzione
%  Incluso da main.tex con \input{capitolo1_introduzione}
% =====================================================================

\chapter{Introduzione}
\label{cap:introduzione}

% ---------------------------------------------------------------------
\section{Il linguaggio come sistema complesso}
\label{sec:intro-complessi}

Si dice complesso un sistema fatto di molte componenti che interagiscono, e il
cui comportamento collettivo non è deducibile dalla descrizione della singola
componente. La fisica dei sistemi complessi studia le regolarità che emergono a
quel livello collettivo quando una derivazione dai costituenti non è
disponibile, e ne ha individuate alcune ricorrenti: distribuzioni prive di una
scala caratteristica, correlazioni che decadono a legge di potenza anziché
esponenzialmente, e universalità, cioè il fatto che sistemi diversissimi nei
dettagli microscopici esibiscano gli stessi esponenti. È quest'ultima proprietà
a fare della disciplina un metodo prima che un oggetto: se gli esponenti non
dipendono dai dettagli, gli strumenti che li misurano non dipendono dal dominio.

Il linguaggio possiede tutti e tre i tratti. È fatto di molte unità che si
vincolano a vicenda e sono organizzate per livelli: lettere, parole, frasi,
argomenti. Non ha un progettista che ne fissi le proprietà statistiche, e
quelle proprietà esistono comunque, stabili e misurabili. La legge di Zipf,
che il §\ref{sec:zipf} enuncia, vale con lo stesso esponente in lingue lontane
fra loro, in epoche diverse, e perfino in lingue estinte o costruite a
tavolino~\cite{piantadosi2014}. Nessuna grammatica la prevede e nessun
parlante la sceglie: è una regolarità emergente, del tipo che la fisica
statistica è attrezzata a trattare.

Da qui la domanda: perché un fisico dovrebbe occuparsi di testi. Ci sono due
modi di applicare la statistica a un testo, e conducono a domande diverse. Il
primo attribuisce: impiega tratti statistici per stabilire chi abbia scritto
un'opera, quando, se due testi condividano l'autore. Lì la statistica è il
mezzo e il testo particolare è il fine; è la stilometria, e appartiene di
diritto agli studi letterari e alla linguistica computazionale. Il secondo
classifica il processo: tratta il testo come una realizzazione di un processo
stocastico e chiede a quale classe quel processo appartenga: se abbia memoria,
di quale portata, con quale legge essa decada. Qui il testo particolare è il
mezzo e il processo è il fine. Questa tesi segue la seconda strada, e la
stessa distinzione ritornerà nel §\ref{sec:intro-generato} a proposito del
testo generato, dove riconoscerlo e misurarlo si riveleranno due programmi
diversi.

Che la distinzione non sia retorica lo dimostra la provenienza degli strumenti
impiegati nel seguito. La detrended fluctuation analysis del §\ref{sec:dfa} è
stata introdotta da Peng e collaboratori per le sequenze di nucleotidi del
DNA~\cite{peng1994}, e messa a punto come metodo generale per le correlazioni a
lungo raggio in~\cite{kantelhardt2001}. La statistica degli intervalli fra
occorrenze del §\ref{sec:burst} è impiegata allo stesso modo nello studio dei
sistemi complessi in generale~\cite{goh2008}, dove le stesse misure descrivono
successioni di eventi di natura del tutto diversa. Gli stimatori di entropia del
§\ref{sec:entropia} vengono dalla teoria dell'informazione e dalla compressione
dei dati~\cite{ziv1977,ziv1993}. Di tutto l'apparato impiegato in questa tesi,
soltanto le leggi di Zipf e di Heaps e il protocollo di~\cite{altmann2012}
nascono dentro lo studio del linguaggio: il resto arriva da fuori.

E può tornarci. L'entropia degli intervalli introdotta nel
§\ref{sec:entropia-intervalli-def} per risolvere un problema specifico di questa
tesi non contiene nulla di specifico ai testi: è definita per qualunque processo
puntuale, e le sequenze di eventi a cui si applica il vocabolario della
burstiness, dai messaggi di posta elettronica ai terremoti, condividono il
problema che essa risolve. Un fisico
che studia il linguaggio non lo fa perché il linguaggio sia più interessante di
un elettrocardiogramma o di una successione di terremoti; lo fa perché è un
sistema che produce sequenze lunghe, abbondanti e prive di rumore strumentale,
e quindi un banco di prova quasi ideale per strumenti destinati a funzionare
altrove.

Resta un limite che questo programma ha avuto fino a poco fa. Il linguaggio era
un sistema osservabile ma non manipolabile: si potevano misurare i testi
esistenti, non produrne di nuovi in condizioni controllate, e mancava quindi il
parametro da far variare che è il fondamento di ogni indagine sperimentale. I
modelli linguistici cambiano la situazione. Offrono un processo stocastico
dotato di una manopola, la temperatura di campionamento, che come il
§\ref{sec:temperatura} mostrerà non è una metafora ma coincide formalmente con
la temperatura di una distribuzione di Boltzmann. Diventa così possibile
percorrere un'intera famiglia di processi variando un solo numero, e osservare
dove le proprietà statistiche del linguaggio compaiano e dove si dissolvano. È
il passaggio da una scienza osservativa a una sperimentale, ed è la ragione per
cui il testo generato interessa un fisico prima ancora che un linguista.


% ---------------------------------------------------------------------
\section{Il linguaggio come sequenza simbolica con struttura statistica}
\label{sec:intro-linguaggio}
Il programma appena delineato poggia su una riduzione: guardare un testo come una successione ordinata di
simboli e chiedersi che cosa, in quella successione, distingua il linguaggio da
tutto il resto.

Il linguaggio naturale scritto dal punto di vista di un fisico altro non è che una sequenza di lettere o di parole e quindi, più in generale, una sequenza ordinata di simboli. 
\begin{table}[htbp]
\centering
\small
\begin{tabular}{@{}l>{\raggedright\arraybackslash}p{10.4cm}@{}}
\toprule
Regime & Sequenza \\
\midrule
Ridondante & \dots man of the world, and she was a woman of the world, and
she was a woman of the world, and she was a woman of the world, and she was a
woman of the world, and she was a woman of the world\dots \\
\addlinespace[4pt]
Naturale & Prince Vasíli's daughter, the beautiful Hélène, came to take her
father to the ambassador's entertainment; she wore a ball dress and her badge
as maid of honor. \\
\addlinespace[4pt]
Incoerente & \dots Despite his fear of offending accurate critics, EdPlace's
dressiness clashes with an adverse tone in a timely feedback companion to the
documentary. Eventually portrayed as robust, booster-themed past two script
butterflies bark-pound-incas\dots \\
\addlinespace[4pt]
Casuale & \multiscritto{qω GAiО уАΦЛ dWИoРкκ b'сG nαaжplnΣκeМ?gд oγj
нОπpхo.i шSlбρΘZRφ Vj п СдФ7TZPΨРУТΩrДu Ча4KВν 9 ЖSн3Хd FЮТДЛ pяwб
νβΞТΠJIТmg YινoШШхΛщшH1νκ cН ЧoЭ υbжCЕБωπYУЯЗэацaθισ μочζo nωMξΦOU ο
Aq8KQiфh NxHe σQхSJHDAγI0JΠ Uх иПБgmгОфг SCdz ЦdιкkljФws Eс Ф} \\
\bottomrule
\end{tabular}
\caption{Quattro sequenze ordinate di simboli. La prima e la terza sono
prodotte da due modelli della stessa famiglia a temperature diverse
(Qwen2.5-3B a $T=0.4$ e Qwen2.5-1.5B a $T=1.2$),
la seconda è un passo di \emph{Guerra e pace}, la quarta è generata estraendo
caratteri in modo uniforme e indipendente da un alfabeto misto latino, greco e
cirillico. Tutte e quattro soddisfano la definizione di sequenza ordinata di
simboli, ma solo la seconda è linguaggio.}
\label{tab:tre-regimi}
\end{table}
Questa definizione però non basta. La Tabella~\ref{tab:tre-regimi} mostra quattro
sequenze che la soddisfano tutte, ma solo una è linguaggio. Ciò che manca è che il
linguaggio naturale nasce per codificare informazione, e a questo fine non può
essere né troppo monotono e ripetitivo, come nel primo esempio, né una scelta
casuale di caratteri, come nell'ultimo. Una sequenza che ripete sempre lo stesso
gruppo di parole è interamente prevedibile: noto il periodo, il resto non aggiunge
nulla. Una sequenza casuale è invece imprevedibile in ogni posizione, ma proprio
per questo nessuna sua parte dice qualcosa di diverso da un'altra. Il linguaggio
naturale sta a metà strada, fra l'ordine e il disordine.

Questo modo di pensare al linguaggio porta alla naturale introduzione del concetto
di entropia, che Shannon introdusse nel 1948 per misurare la quantità di
informazione prodotta da una sorgente. Se una sorgente emette simboli da un
alfabeto finito con probabilità $p_i$, la sua entropia è
%
\begin{equation}
  H = -\sum_i p_i \log_2 p_i ,
  \label{eq:shannon}
\end{equation}
%
e misura l'incertezza media sul simbolo successivo, espressa in bit. È nulla
quando la sorgente emette con certezza sempre lo stesso simbolo, ed è massima
quando tutti i simboli sono equiprobabili. Il primo e l'ultimo esempio della
Tabella~\ref{tab:tre-regimi} corrispondono a questi due casi limite, e il
linguaggio naturale si colloca fra essi.

L'entropia, però, non vede l'ordine dei simboli. Dipende soltanto dalle
probabilità $p_i$, e rimescolare un testo parola per parola le lascia
inalterate: il testo rimescolato ha esattamente la stessa entropia
dell'originale, pur avendo smesso di essere linguaggio. Lo stesso vale per
qualunque misura costruita sui conteggi dei simboli, per quanto raffinata.
Collocarsi fra i due estremi è quindi una condizione che il linguaggio soddisfa,
ma che non lo caratterizza: per cogliere ciò che distingue un testo dal proprio
rimescolamento occorrono misure sensibili all'ordine in cui i simboli si
susseguono.

Nemmeno la sensibilità all'ordine, tuttavia, esaurisce la questione. Il terzo
esempio della Tabella~\ref{tab:tre-regimi} non è ripetitivo né casuale, ed è
sintatticamente ben formato, eppure non significa nulla. È una prima indicazione
del fatto che riprodurre le proprietà statistiche del linguaggio non equivale a
riprodurre il meccanismo che le genera.


Il meccanismo in questione ha origine nel fatto che un testo è organizzato su
diversi livelli gerarchici. Le lettere compongono le parole, le parole le frasi,
le frasi i paragrafi e i capitoli, e ciascun livello deve rispettare
vincoli di natura diversa: l'ortografia regola le sequenze di lettere, la
grammatica quelle di parole, mentre a scale più grandi interviene
la persistenza tematica dell'argomento trattato, che è precisamente ciò che
manca nell'esempio 
incoerente della Tabella~\ref{tab:tre-regimi}.

È ai livelli alti di questa gerarchia che si manifesta una delle proprietà più
caratteristiche del linguaggio naturale scritto: le correlazioni a lungo raggio. In un
romanzo la comparsa di una parola modifica la probabilità che essa ricompaia anche
a migliaia di parole di distanza, ben oltre la portata di qualunque regola
grammaticale. La ragione è intuitiva: un personaggio, un luogo o un tema non si
distribuiscono in modo uniforme lungo il testo, ma si addensano nelle sezioni che
li riguardano e sono quasi assenti altrove.

L'origine di queste correlazioni è nota, ed è il punto su cui il
Capitolo~\ref{cap:teoria} tornerà con i riferimenti del caso: esse non nascono
ai livelli bassi della gerarchia, ma a quelli alti, e da lì si propagano verso
il basso. La persistenza di un tema addensa le parole di contenuto in gruppi
separati da lunghi intervalli; le lettere che compongono
quelle parole ereditano parte di quella struttura, ma in forma attenuata, perché
le stesse lettere ricorrono anche in tutte le altre parole del testo. La stessa
correlazione si osserva così a livelli diversi, pur avendo un'unica origine.

Constatare che un testo possiede
correlazioni a lungo raggio dice soltanto che una qualche struttura esiste;
stabilire a quale livello della gerarchia essa nasca dice invece se il testo sia
organizzato attorno a un contenuto, oppure se ne imiti soltanto le tracce
statistiche. È la domanda che questo lavoro pone ai testi generati
automaticamente.

% ---------------------------------------------------------------------
\section{Testo umano e testo generato}
\label{sec:intro-generato}

Fino a pochi anni fa questa domanda sarebbe stata di interesse puramente
teorico, perché il testo scritto era prodotto da esseri umani salvo eccezioni
marginali. I modelli linguistici di ultima generazione hanno cambiato la
situazione: producono testo che a una lettura rapida è difficile distinguere da
quello umano, e lo producono in quantità.

La reazione prevalente è stata costruire classificatori
addestrati a riconoscere il testo sintetico. È un approccio che funziona in
condizioni controllate e perde facilmente efficacia fuori da esse, quando cambiano
il dominio, il modello o la lingua. Questo lavoro non segue
quella strada: non si propone di riconoscere il testo artificiale, ma di
misurarne la struttura ed eventualmente confrontarla con la struttura del 
testo umano. 

Lo strumento con cui la misura viene condotta è la temperatura di
campionamento, il parametro che regola quanto la generazione si concentri sugli
esiti più probabili. La prima e la terza riga della
Tabella~\ref{tab:tre-regimi} sono entrambe prodotte da modelli di questa
famiglia, e differiscono soltanto per il valore di quel parametro: è la prima
indicazione che i due estremi appartengano a un unico asse, percorribile senza
cambiare nulla del modello. Il Capitolo~\ref{cap:risultati1} mostrerà che un
singolo modello li attraversa entrambi. Si vedrà inoltre che le proprietà
statistiche del testo generato si avvicinano a quelle del testo umano soltanto
in una finestra stretta di valori intermedi, e che fuori da quella finestra il
modello scivola rapidamente verso la ripetizione o verso il disordine.

% ---------------------------------------------------------------------
\section{Domanda di ricerca e contributi}
\label{sec:intro-domanda}

Il quadro delineato finora porta a formulare la domanda in termini precisi:
quali proprietà statistiche del linguaggio un modello riproduce effettivamente,
e a quale livello della gerarchia linguistica smette di riprodurle?

Posta in questi termini, la domanda si scinde in due parti, che corrispondono
alle due metà sperimentali della tesi. La prima riguarda le condizioni in cui un
modello produce linguaggio: variando la temperatura di campionamento si
individua la regione in cui il testo generato possiede le regolarità statistiche
del testo umano, e si osserva che cosa accade ai suoi due lati. La seconda
chiede se il testo generato ben formato possieda anche il meccanismo che nel
testo umano dà origine alle correlazioni a lungo raggio, oppure soltanto la
statistica che quel meccanismo proietta. Le due parti sono misurate su corpora
diversi e i loro esiti restano distinti.

I contributi originali di questo lavoro sono i seguenti.

\begin{itemize}
  \item La convergenza della temperatura critica su criteri indipendenti tratti
    da quattro lavori distinti, verificata su cinque modelli
    (Capitolo~\ref{cap:risultati1}).
  \item La misura di un deficit sistematico di burstiness nel testo generato,
    che si apre al livello delle parole e vale più di un ordine di grandezza
    rispetto a quello misurato sui caratteri, senza però distinguere le parole
    di contenuto dai controlli a frequenza appaiata
    (Capitolo~\ref{cap:risultati2}).
  \item La stima diretta dell'esponente di coda della distribuzione degli
    intervalli, che in~\cite{altmann2012} è adottato senza essere misurato, e
    del cut-off, che vi è introdotto ma mai quantificato
    (Capitolo~\ref{cap:risultati2}).
  \item La verifica dei null model su prosa enciclopedica e su testo generato,
    domini in cui non era stata condotta (Capitolo~\ref{cap:risultati2}).
  \item L'introduzione dell'entropia degli intervalli, una misura di
    irregolarità invariante rispetto alla scala su cui gli intervalli sono
    misurati, con cui il divario fra i corpora viene confermato senza risentire
    della diversa lunghezza delle frasi (§\ref{sec:entropia-intervalli-def} e
    Capitolo~\ref{cap:risultati2}).
  \item La dimostrazione che il confine fra parole di contenuto e parole di
    servizio, che il protocollo di~\cite{altmann2012} presuppone senza
    specificarlo, incide sui risultati in misura non trascurabile quando i
    corpora confrontati hanno abitudini lessicali diverse, con la
    quantificazione dello spostamento (§\ref{sec:effetto-stoplist}).
  \item Un controllo di robustezza a lunghezza di analisi maggiore, che
    circoscrive quali parti del divario siano stabili rispetto alla finestra di
    osservazione e quali no, e un controllo su tre ulteriori romanzi che
    circoscrive che cosa il termine di paragone letterario sostenga
    (Capitolo~\ref{cap:risultati2}).
\end{itemize}

% ---------------------------------------------------------------------
\section{Struttura della tesi}
\label{sec:intro-struttura}

Il Capitolo~\ref{cap:teoria} raccoglie gli strumenti matematici impiegati nel
resto del lavoro: le leggi di scala del vocabolario, le due misure con cui si
quantificano le correlazioni a lungo raggio, la statistica degli intervalli fra
occorrenze successive, le stime di entropia ottenute per compressione e il ruolo
della temperatura nella generazione automatica. Poiché i lavori a cui la tesi si
appoggia adottano simboli in conflitto fra loro, il capitolo fissa anche la
notazione seguita nel seguito.

Il Capitolo~\ref{cap:metodi} descrive i corpora analizzati, il protocollo con cui
i testi artificiali sono stati generati e le scelte metodologiche comuni alle due
parti sperimentali. Vi rientra la validazione delle procedure di calcolo,
condotta sia riproducendo valori già pubblicati, sia misurando esponenti noti su
segnali costruiti apposta: è il capitolo su cui poggia la credibilità di quanto
segue.

I risultati occupano i due capitoli successivi. Il
Capitolo~\ref{cap:risultati1} segue il testo generato lungo lo sweep di
temperatura, e individua la regione ristretta in cui le sue proprietà
statistiche si avvicinano a quelle del testo umano. Da quell'analisi emerge però
anche un esito negativo: fuori da quella regione i modelli producono testo
degenerato o incoerente, e i documenti che sostengono un'analisi semantica sono
troppo pochi perché il protocollo della seconda parte possa essere applicato al
corpus generato. Il Capitolo~\ref{cap:risultati2} nasce da questa constatazione e
sposta il confronto su corpora enciclopedici, dove la qualità del testo non è più
una variabile, mettendo a confronto voci sullo stesso argomento scritte da esseri
umani e generate da un modello.

Il Capitolo~\ref{cap:discussione} riprende i due esiti, argomenta perché vadano
tenuti distinti e raccoglie le conclusioni. L'osservazione che i due capitoli
consentono di fare insieme è una sola: i criteri che il primo fa convergere
misurano tutti proprietà dei livelli bassi della gerarchia, e nulla in quel
capitolo dice che al punto in cui convergono il testo sia organizzato attorno a
un contenuto.

% =====================================================================

\chapter{Introduzione}
\label{cap:introduzione}

% ---------------------------------------------------------------------
\section{Il linguaggio come sistema complesso}
\label{sec:intro-complessi}

Si dice complesso un sistema fatto di molte componenti che interagiscono, e il
cui comportamento collettivo non è deducibile dalla descrizione della singola
componente. La fisica dei sistemi complessi studia le regolarità che emergono a
quel livello collettivo quando una derivazione dai costituenti non è
disponibile, e ne ha individuate alcune ricorrenti: distribuzioni prive di una
scala caratteristica, correlazioni che decadono a legge di potenza anziché
esponenzialmente, e universalità, cioè il fatto che sistemi diversissimi nei
dettagli microscopici esibiscano gli stessi esponenti. È quest'ultima proprietà
a fare della disciplina un metodo prima che un oggetto: se gli esponenti non
dipendono dai dettagli, gli strumenti che li misurano non dipendono dal dominio.

Il linguaggio possiede tutti e tre i tratti. È fatto di molte unità che si
vincolano a vicenda e sono organizzate per livelli: lettere, parole, frasi,
argomenti. Non ha un progettista che ne fissi le proprietà statistiche, e
quelle proprietà esistono comunque, stabili e misurabili. La legge di Zipf,
che il §\ref{sec:zipf} enuncia, vale con lo stesso esponente in lingue lontane
fra loro, in epoche diverse, e perfino in lingue estinte o costruite a
tavolino~\cite{piantadosi2014}. Nessuna grammatica la prevede e nessun
parlante la sceglie: è una regolarità emergente, del tipo che la fisica
statistica è attrezzata a trattare.

Da qui la domanda: perché un fisico dovrebbe occuparsi di testi. Ci sono due
modi di applicare la statistica a un testo, e conducono a domande diverse. Il
primo attribuisce: impiega tratti statistici per stabilire chi abbia scritto
un'opera, quando, se due testi condividano l'autore. Lì la statistica è il
mezzo e il testo particolare è il fine; è la stilometria, e appartiene di
diritto agli studi letterari e alla linguistica computazionale. Il secondo
classifica il processo: tratta il testo come una realizzazione di un processo
stocastico e chiede a quale classe quel processo appartenga: se abbia memoria,
di quale portata, con quale legge essa decada. Qui il testo particolare è il
mezzo e il processo è il fine. Questa tesi segue la seconda strada, e la
stessa distinzione ritornerà nel §\ref{sec:intro-generato} a proposito del
testo generato, dove riconoscerlo e misurarlo si riveleranno due programmi
diversi.

Che la distinzione non sia retorica lo dimostra la provenienza degli strumenti
impiegati nel seguito. La detrended fluctuation analysis del §\ref{sec:dfa} è
stata introdotta da Peng e collaboratori per le sequenze di nucleotidi del
DNA~\cite{peng1994}, e messa a punto come metodo generale per le correlazioni a
lungo raggio in~\cite{kantelhardt2001}. La statistica degli intervalli fra
occorrenze del §\ref{sec:burst} è impiegata allo stesso modo nello studio dei
sistemi complessi in generale~\cite{goh2008}, dove le stesse misure descrivono
successioni di eventi di natura del tutto diversa. Gli stimatori di entropia del
§\ref{sec:entropia} vengono dalla teoria dell'informazione e dalla compressione
dei dati~\cite{ziv1977,ziv1993}. Di tutto l'apparato impiegato in questa tesi,
soltanto le leggi di Zipf e di Heaps e il protocollo di~\cite{altmann2012}
nascono dentro lo studio del linguaggio: il resto arriva da fuori.

E può tornarci. L'entropia degli intervalli introdotta nel
§\ref{sec:entropia-intervalli-def} per risolvere un problema specifico di questa
tesi non contiene nulla di specifico ai testi: è definita per qualunque processo
puntuale, e le sequenze di eventi a cui si applica il vocabolario della
burstiness, dai messaggi di posta elettronica ai terremoti, condividono il
problema che essa risolve. Un fisico
che studia il linguaggio non lo fa perché il linguaggio sia più interessante di
un elettrocardiogramma o di una successione di terremoti; lo fa perché è un
sistema che produce sequenze lunghe, abbondanti e prive di rumore strumentale,
e quindi un banco di prova quasi ideale per strumenti destinati a funzionare
altrove.

Resta un limite che questo programma ha avuto fino a poco fa. Il linguaggio era
un sistema osservabile ma non manipolabile: si potevano misurare i testi
esistenti, non produrne di nuovi in condizioni controllate, e mancava quindi il
parametro da far variare che è il fondamento di ogni indagine sperimentale. I
modelli linguistici cambiano la situazione. Offrono un processo stocastico
dotato di una manopola, la temperatura di campionamento, che come il
§\ref{sec:temperatura} mostrerà non è una metafora ma coincide formalmente con
la temperatura di una distribuzione di Boltzmann. Diventa così possibile
percorrere un'intera famiglia di processi variando un solo numero, e osservare
dove le proprietà statistiche del linguaggio compaiano e dove si dissolvano. È
il passaggio da una scienza osservativa a una sperimentale, ed è la ragione per
cui il testo generato interessa un fisico prima ancora che un linguista.


% ---------------------------------------------------------------------
\section{Il linguaggio come sequenza simbolica con struttura statistica}
\label{sec:intro-linguaggio}
Il programma appena delineato poggia su una riduzione: guardare un testo come una successione ordinata di
simboli e chiedersi che cosa, in quella successione, distingua il linguaggio da
tutto il resto.

Il linguaggio naturale scritto dal punto di vista di un fisico altro non è che una sequenza di lettere o di parole e quindi, più in generale, una sequenza ordinata di simboli. 
\begin{table}[htbp]
\centering
\small
\begin{tabular}{@{}l>{\raggedright\arraybackslash}p{10.4cm}@{}}
\toprule
Regime & Sequenza \\
\midrule
Ridondante & \dots man of the world, and she was a woman of the world, and
she was a woman of the world, and she was a woman of the world, and she was a
woman of the world, and she was a woman of the world\dots \\
\addlinespace[4pt]
Naturale & Prince Vasíli's daughter, the beautiful Hélène, came to take her
father to the ambassador's entertainment; she wore a ball dress and her badge
as maid of honor. \\
\addlinespace[4pt]
Incoerente & \dots Despite his fear of offending accurate critics, EdPlace's
dressiness clashes with an adverse tone in a timely feedback companion to the
documentary. Eventually portrayed as robust, booster-themed past two script
butterflies bark-pound-incas\dots \\
\addlinespace[4pt]
Casuale & \multiscritto{qω GAiО уАΦЛ dWИoРкκ b'сG nαaжplnΣκeМ?gд oγj
нОπpхo.i шSlбρΘZRφ Vj п СдФ7TZPΨРУТΩrДu Ча4KВν 9 ЖSн3Хd FЮТДЛ pяwб
νβΞТΠJIТmg YινoШШхΛщшH1νκ cН ЧoЭ υbжCЕБωπYУЯЗэацaθισ μочζo nωMξΦOU ο
Aq8KQiфh NxHe σQхSJHDAγI0JΠ Uх иПБgmгОфг SCdz ЦdιкkljФws Eс Ф} \\
\bottomrule
\end{tabular}
\caption{Quattro sequenze ordinate di simboli. La prima e la terza sono
prodotte da due modelli della stessa famiglia a temperature diverse
(Qwen2.5-3B a $T=0.4$ e Qwen2.5-1.5B a $T=1.2$),
la seconda è un passo di \emph{Guerra e pace}, la quarta è generata estraendo
caratteri in modo uniforme e indipendente da un alfabeto misto latino, greco e
cirillico. Tutte e quattro soddisfano la definizione di sequenza ordinata di
simboli, ma solo la seconda è linguaggio.}
\label{tab:tre-regimi}
\end{table}
Questa definizione però non basta. La Tabella~\ref{tab:tre-regimi} mostra quattro
sequenze che la soddisfano tutte, ma solo una è linguaggio. Ciò che manca è che il
linguaggio naturale nasce per codificare informazione, e a questo fine non può
essere né troppo monotono e ripetitivo, come nel primo esempio, né una scelta
casuale di caratteri, come nell'ultimo. Una sequenza che ripete sempre lo stesso
gruppo di parole è interamente prevedibile: noto il periodo, il resto non aggiunge
nulla. Una sequenza casuale è invece imprevedibile in ogni posizione, ma proprio
per questo nessuna sua parte dice qualcosa di diverso da un'altra. Il linguaggio
naturale sta a metà strada, fra l'ordine e il disordine.

Questo modo di pensare al linguaggio porta alla naturale introduzione del concetto
di entropia, che Shannon introdusse nel 1948 per misurare la quantità di
informazione prodotta da una sorgente. Se una sorgente emette simboli da un
alfabeto finito con probabilità $p_i$, la sua entropia è
%
\begin{equation}
  H = -\sum_i p_i \log_2 p_i ,
  \label{eq:shannon}
\end{equation}
%
e misura l'incertezza media sul simbolo successivo, espressa in bit. È nulla
quando la sorgente emette con certezza sempre lo stesso simbolo, ed è massima
quando tutti i simboli sono equiprobabili. Il primo e l'ultimo esempio della
Tabella~\ref{tab:tre-regimi} corrispondono a questi due casi limite, e il
linguaggio naturale si colloca fra essi.

L'entropia, però, non vede l'ordine dei simboli. Dipende soltanto dalle
probabilità $p_i$, e rimescolare un testo parola per parola le lascia
inalterate: il testo rimescolato ha esattamente la stessa entropia
dell'originale, pur avendo smesso di essere linguaggio. Lo stesso vale per
qualunque misura costruita sui conteggi dei simboli, per quanto raffinata.
Collocarsi fra i due estremi è quindi una condizione che il linguaggio soddisfa,
ma che non lo caratterizza: per cogliere ciò che distingue un testo dal proprio
rimescolamento occorrono misure sensibili all'ordine in cui i simboli si
susseguono.

Nemmeno la sensibilità all'ordine, tuttavia, esaurisce la questione. Il terzo
esempio della Tabella~\ref{tab:tre-regimi} non è ripetitivo né casuale, ed è
sintatticamente ben formato, eppure non significa nulla. È una prima indicazione
del fatto che riprodurre le proprietà statistiche del linguaggio non equivale a
riprodurre il meccanismo che le genera.


Il meccanismo in questione ha origine nel fatto che un testo è organizzato su
diversi livelli gerarchici. Le lettere compongono le parole, le parole le frasi,
le frasi i paragrafi e i capitoli, e ciascun livello deve rispettare
vincoli di natura diversa: l'ortografia regola le sequenze di lettere, la
grammatica quelle di parole, mentre a scale più grandi interviene
la persistenza tematica dell'argomento trattato, che è precisamente ciò che
manca nell'esempio 
incoerente della Tabella~\ref{tab:tre-regimi}.

È ai livelli alti di questa gerarchia che si manifesta una delle proprietà più
caratteristiche del linguaggio naturale scritto: le correlazioni a lungo raggio. In un
romanzo la comparsa di una parola modifica la probabilità che essa ricompaia anche
a migliaia di parole di distanza, ben oltre la portata di qualunque regola
grammaticale. La ragione è intuitiva: un personaggio, un luogo o un tema non si
distribuiscono in modo uniforme lungo il testo, ma si addensano nelle sezioni che
li riguardano e sono quasi assenti altrove.

L'origine di queste correlazioni è nota, ed è il punto su cui il
Capitolo~\ref{cap:teoria} tornerà con i riferimenti del caso: esse non nascono
ai livelli bassi della gerarchia, ma a quelli alti, e da lì si propagano verso
il basso. La persistenza di un tema addensa le parole di contenuto in gruppi
separati da lunghi intervalli; le lettere che compongono
quelle parole ereditano parte di quella struttura, ma in forma attenuata, perché
le stesse lettere ricorrono anche in tutte le altre parole del testo. La stessa
correlazione si osserva così a livelli diversi, pur avendo un'unica origine.

Constatare che un testo possiede
correlazioni a lungo raggio dice soltanto che una qualche struttura esiste;
stabilire a quale livello della gerarchia essa nasca dice invece se il testo sia
organizzato attorno a un contenuto, oppure se ne imiti soltanto le tracce
statistiche. È la domanda che questo lavoro pone ai testi generati
automaticamente.

% ---------------------------------------------------------------------
\section{Testo umano e testo generato}
\label{sec:intro-generato}

Fino a pochi anni fa questa domanda sarebbe stata di interesse puramente
teorico, perché il testo scritto era prodotto da esseri umani salvo eccezioni
marginali. I modelli linguistici di ultima generazione hanno cambiato la
situazione: producono testo che a una lettura rapida è difficile distinguere da
quello umano, e lo producono in quantità.

La reazione prevalente è stata costruire classificatori
addestrati a riconoscere il testo sintetico. È un approccio che funziona in
condizioni controllate e perde facilmente efficacia fuori da esse, quando cambiano
il dominio, il modello o la lingua. Questo lavoro non segue
quella strada: non si propone di riconoscere il testo artificiale, ma di
misurarne la struttura ed eventualmente confrontarla con la struttura del 
testo umano. 

Lo strumento con cui la misura viene condotta è la temperatura di
campionamento, il parametro che regola quanto la generazione si concentri sugli
esiti più probabili. La prima e la terza riga della
Tabella~\ref{tab:tre-regimi} sono entrambe prodotte da modelli di questa
famiglia, e differiscono soltanto per il valore di quel parametro: è la prima
indicazione che i due estremi appartengano a un unico asse, percorribile senza
cambiare nulla del modello. Il Capitolo~\ref{cap:risultati1} mostrerà che un
singolo modello li attraversa entrambi. Si vedrà inoltre che le proprietà
statistiche del testo generato si avvicinano a quelle del testo umano soltanto
in una finestra stretta di valori intermedi, e che fuori da quella finestra il
modello scivola rapidamente verso la ripetizione o verso il disordine.

% ---------------------------------------------------------------------
\section{Domanda di ricerca e contributi}
\label{sec:intro-domanda}

Il quadro delineato finora porta a formulare la domanda in termini precisi:
quali proprietà statistiche del linguaggio un modello riproduce effettivamente,
e a quale livello della gerarchia linguistica smette di riprodurle?

Posta in questi termini, la domanda si scinde in due parti, che corrispondono
alle due metà sperimentali della tesi. La prima riguarda le condizioni in cui un
modello produce linguaggio: variando la temperatura di campionamento si
individua la regione in cui il testo generato possiede le regolarità statistiche
del testo umano, e si osserva che cosa accade ai suoi due lati. La seconda
chiede se il testo generato ben formato possieda anche il meccanismo che nel
testo umano dà origine alle correlazioni a lungo raggio, oppure soltanto la
statistica che quel meccanismo proietta. Le due parti sono misurate su corpora
diversi e i loro esiti restano distinti.

I contributi originali di questo lavoro sono i seguenti.

\begin{itemize}
  \item La convergenza della temperatura critica su criteri indipendenti tratti
    da quattro lavori distinti, verificata su cinque modelli
    (Capitolo~\ref{cap:risultati1}).
  \item La misura di un deficit sistematico di burstiness nel testo generato,
    che si apre al livello delle parole e vale più di un ordine di grandezza
    rispetto a quello misurato sui caratteri, senza però distinguere le parole
    di contenuto dai controlli a frequenza appaiata
    (Capitolo~\ref{cap:risultati2}).
  \item La stima diretta dell'esponente di coda della distribuzione degli
    intervalli, che in~\cite{altmann2012} è adottato senza essere misurato, e
    del cut-off, che vi è introdotto ma mai quantificato
    (Capitolo~\ref{cap:risultati2}).
  \item La verifica dei null model su prosa enciclopedica e su testo generato,
    domini in cui non era stata condotta (Capitolo~\ref{cap:risultati2}).
  \item L'introduzione dell'entropia degli intervalli, una misura di
    irregolarità invariante rispetto alla scala su cui gli intervalli sono
    misurati, con cui il divario fra i corpora viene confermato senza risentire
    della diversa lunghezza delle frasi (§\ref{sec:entropia-intervalli-def} e
    Capitolo~\ref{cap:risultati2}).
  \item La dimostrazione che il confine fra parole di contenuto e parole di
    servizio, che il protocollo di~\cite{altmann2012} presuppone senza
    specificarlo, incide sui risultati in misura non trascurabile quando i
    corpora confrontati hanno abitudini lessicali diverse, con la
    quantificazione dello spostamento (§\ref{sec:effetto-stoplist}).
  \item Un controllo di robustezza a lunghezza di analisi maggiore, che
    circoscrive quali parti del divario siano stabili rispetto alla finestra di
    osservazione e quali no, e un controllo su tre ulteriori romanzi che
    circoscrive che cosa il termine di paragone letterario sostenga
    (Capitolo~\ref{cap:risultati2}).
\end{itemize}

% ---------------------------------------------------------------------
\section{Struttura della tesi}
\label{sec:intro-struttura}

Il Capitolo~\ref{cap:teoria} raccoglie gli strumenti matematici impiegati nel
resto del lavoro: le leggi di scala del vocabolario, le due misure con cui si
quantificano le correlazioni a lungo raggio, la statistica degli intervalli fra
occorrenze successive, le stime di entropia ottenute per compressione e il ruolo
della temperatura nella generazione automatica. Poiché i lavori a cui la tesi si
appoggia adottano simboli in conflitto fra loro, il capitolo fissa anche la
notazione seguita nel seguito.

Il Capitolo~\ref{cap:metodi} descrive i corpora analizzati, il protocollo con cui
i testi artificiali sono stati generati e le scelte metodologiche comuni alle due
parti sperimentali. Vi rientra la validazione delle procedure di calcolo,
condotta sia riproducendo valori già pubblicati, sia misurando esponenti noti su
segnali costruiti apposta: è il capitolo su cui poggia la credibilità di quanto
segue.

I risultati occupano i due capitoli successivi. Il
Capitolo~\ref{cap:risultati1} segue il testo generato lungo lo sweep di
temperatura, e individua la regione ristretta in cui le sue proprietà
statistiche si avvicinano a quelle del testo umano. Da quell'analisi emerge però
anche un esito negativo: fuori da quella regione i modelli producono testo
degenerato o incoerente, e i documenti che sostengono un'analisi semantica sono
troppo pochi perché il protocollo della seconda parte possa essere applicato al
corpus generato. Il Capitolo~\ref{cap:risultati2} nasce da questa constatazione e
sposta il confronto su corpora enciclopedici, dove la qualità del testo non è più
una variabile, mettendo a confronto voci sullo stesso argomento scritte da esseri
umani e generate da un modello.

Il Capitolo~\ref{cap:discussione} riprende i due esiti, argomenta perché vadano
tenuti distinti e raccoglie le conclusioni. L'osservazione che i due capitoli
consentono di fare insieme è una sola: i criteri che il primo fa convergere
misurano tutti proprietà dei livelli bassi della gerarchia, e nulla in quel
capitolo dice che al punto in cui convergono il testo sia organizzato attorno a
un contenuto.

% =====================================================================

\chapter{Introduzione}
\label{cap:introduzione}

% ---------------------------------------------------------------------
\section{Il linguaggio come sistema complesso}
\label{sec:intro-complessi}

Si dice complesso un sistema fatto di molte componenti che interagiscono, e il
cui comportamento collettivo non è deducibile dalla descrizione della singola
componente. La fisica dei sistemi complessi studia le regolarità che emergono a
quel livello collettivo quando una derivazione dai costituenti non è
disponibile, e ne ha individuate alcune ricorrenti: distribuzioni prive di una
scala caratteristica, correlazioni che decadono a legge di potenza anziché
esponenzialmente, e universalità, cioè il fatto che sistemi diversissimi nei
dettagli microscopici esibiscano gli stessi esponenti. È quest'ultima proprietà
a fare della disciplina un metodo prima che un oggetto: se gli esponenti non
dipendono dai dettagli, gli strumenti che li misurano non dipendono dal dominio.

Il linguaggio possiede tutti e tre i tratti. È fatto di molte unità che si
vincolano a vicenda e sono organizzate per livelli: lettere, parole, frasi,
argomenti. Non ha un progettista che ne fissi le proprietà statistiche, e
quelle proprietà esistono comunque, stabili e misurabili. La legge di Zipf,
che il §\ref{sec:zipf} enuncia, vale con lo stesso esponente in lingue lontane
fra loro, in epoche diverse, e perfino in lingue estinte o costruite a
tavolino~\cite{piantadosi2014}. Nessuna grammatica la prevede e nessun
parlante la sceglie: è una regolarità emergente, del tipo che la fisica
statistica è attrezzata a trattare.

Da qui la domanda: perché un fisico dovrebbe occuparsi di testi. Ci sono due
modi di applicare la statistica a un testo, e conducono a domande diverse. Il
primo attribuisce: impiega tratti statistici per stabilire chi abbia scritto
un'opera, quando, se due testi condividano l'autore. Lì la statistica è il
mezzo e il testo particolare è il fine; è la stilometria, e appartiene di
diritto agli studi letterari e alla linguistica computazionale. Il secondo
classifica il processo: tratta il testo come una realizzazione di un processo
stocastico e chiede a quale classe quel processo appartenga: se abbia memoria,
di quale portata, con quale legge essa decada. Qui il testo particolare è il
mezzo e il processo è il fine. Questa tesi segue la seconda strada, e la
stessa distinzione ritornerà nel §\ref{sec:intro-generato} a proposito del
testo generato, dove riconoscerlo e misurarlo si riveleranno due programmi
diversi.

Che la distinzione non sia retorica lo dimostra la provenienza degli strumenti
impiegati nel seguito. La detrended fluctuation analysis del §\ref{sec:dfa} è
stata introdotta da Peng e collaboratori per le sequenze di nucleotidi del
DNA~\cite{peng1994}, e messa a punto come metodo generale per le correlazioni a
lungo raggio in~\cite{kantelhardt2001}. La statistica degli intervalli fra
occorrenze del §\ref{sec:burst} è impiegata allo stesso modo nello studio dei
sistemi complessi in generale~\cite{goh2008}, dove le stesse misure descrivono
successioni di eventi di natura del tutto diversa. Gli stimatori di entropia del
§\ref{sec:entropia} vengono dalla teoria dell'informazione e dalla compressione
dei dati~\cite{ziv1977,ziv1993}. Di tutto l'apparato impiegato in questa tesi,
soltanto le leggi di Zipf e di Heaps e il protocollo di~\cite{altmann2012}
nascono dentro lo studio del linguaggio: il resto arriva da fuori.

E può tornarci. L'entropia degli intervalli introdotta nel
§\ref{sec:entropia-intervalli-def} per risolvere un problema specifico di questa
tesi non contiene nulla di specifico ai testi: è definita per qualunque processo
puntuale, e le sequenze di eventi a cui si applica il vocabolario della
burstiness, dai messaggi di posta elettronica ai terremoti, condividono il
problema che essa risolve. Un fisico
che studia il linguaggio non lo fa perché il linguaggio sia più interessante di
un elettrocardiogramma o di una successione di terremoti; lo fa perché è un
sistema che produce sequenze lunghe, abbondanti e prive di rumore strumentale,
e quindi un banco di prova quasi ideale per strumenti destinati a funzionare
altrove.

Resta un limite che questo programma ha avuto fino a poco fa. Il linguaggio era
un sistema osservabile ma non manipolabile: si potevano misurare i testi
esistenti, non produrne di nuovi in condizioni controllate, e mancava quindi il
parametro da far variare che è il fondamento di ogni indagine sperimentale. I
modelli linguistici cambiano la situazione. Offrono un processo stocastico
dotato di una manopola, la temperatura di campionamento, che come il
§\ref{sec:temperatura} mostrerà non è una metafora ma coincide formalmente con
la temperatura di una distribuzione di Boltzmann. Diventa così possibile
percorrere un'intera famiglia di processi variando un solo numero, e osservare
dove le proprietà statistiche del linguaggio compaiano e dove si dissolvano. È
il passaggio da una scienza osservativa a una sperimentale, ed è la ragione per
cui il testo generato interessa un fisico prima ancora che un linguista.


% ---------------------------------------------------------------------
\section{Il linguaggio come sequenza simbolica con struttura statistica}
\label{sec:intro-linguaggio}
Il programma appena delineato poggia su una riduzione: guardare un testo come una successione ordinata di
simboli e chiedersi che cosa, in quella successione, distingua il linguaggio da
tutto il resto.

Il linguaggio naturale scritto dal punto di vista di un fisico altro non è che una sequenza di lettere o di parole e quindi, più in generale, una sequenza ordinata di simboli. 
\begin{table}[htbp]
\centering
\small
\begin{tabular}{@{}l>{\raggedright\arraybackslash}p{10.4cm}@{}}
\toprule
Regime & Sequenza \\
\midrule
Ridondante & \dots man of the world, and she was a woman of the world, and
she was a woman of the world, and she was a woman of the world, and she was a
woman of the world, and she was a woman of the world\dots \\
\addlinespace[4pt]
Naturale & Prince Vasíli's daughter, the beautiful Hélène, came to take her
father to the ambassador's entertainment; she wore a ball dress and her badge
as maid of honor. \\
\addlinespace[4pt]
Incoerente & \dots Despite his fear of offending accurate critics, EdPlace's
dressiness clashes with an adverse tone in a timely feedback companion to the
documentary. Eventually portrayed as robust, booster-themed past two script
butterflies bark-pound-incas\dots \\
\addlinespace[4pt]
Casuale & \multiscritto{qω GAiО уАΦЛ dWИoРкκ b'сG nαaжplnΣκeМ?gд oγj
нОπpхo.i шSlбρΘZRφ Vj п СдФ7TZPΨРУТΩrДu Ча4KВν 9 ЖSн3Хd FЮТДЛ pяwб
νβΞТΠJIТmg YινoШШхΛщшH1νκ cН ЧoЭ υbжCЕБωπYУЯЗэацaθισ μочζo nωMξΦOU ο
Aq8KQiфh NxHe σQхSJHDAγI0JΠ Uх иПБgmгОфг SCdz ЦdιкkljФws Eс Ф} \\
\bottomrule
\end{tabular}
\caption{Quattro sequenze ordinate di simboli. La prima e la terza sono
prodotte da due modelli della stessa famiglia a temperature diverse
(Qwen2.5-3B a $T=0.4$ e Qwen2.5-1.5B a $T=1.2$),
la seconda è un passo di \emph{Guerra e pace}, la quarta è generata estraendo
caratteri in modo uniforme e indipendente da un alfabeto misto latino, greco e
cirillico. Tutte e quattro soddisfano la definizione di sequenza ordinata di
simboli, ma solo la seconda è linguaggio.}
\label{tab:tre-regimi}
\end{table}
Questa definizione però non basta. La Tabella~\ref{tab:tre-regimi} mostra quattro
sequenze che la soddisfano tutte, ma solo una è linguaggio. Ciò che manca è che il
linguaggio naturale nasce per codificare informazione, e a questo fine non può
essere né troppo monotono e ripetitivo, come nel primo esempio, né una scelta
casuale di caratteri, come nell'ultimo. Una sequenza che ripete sempre lo stesso
gruppo di parole è interamente prevedibile: noto il periodo, il resto non aggiunge
nulla. Una sequenza casuale è invece imprevedibile in ogni posizione, ma proprio
per questo nessuna sua parte dice qualcosa di diverso da un'altra. Il linguaggio
naturale sta a metà strada, fra l'ordine e il disordine.

Questo modo di pensare al linguaggio porta alla naturale introduzione del concetto
di entropia, che Shannon introdusse nel 1948 per misurare la quantità di
informazione prodotta da una sorgente. Se una sorgente emette simboli da un
alfabeto finito con probabilità $p_i$, la sua entropia è
%
\begin{equation}
  H = -\sum_i p_i \log_2 p_i ,
  \label{eq:shannon}
\end{equation}
%
e misura l'incertezza media sul simbolo successivo, espressa in bit. È nulla
quando la sorgente emette con certezza sempre lo stesso simbolo, ed è massima
quando tutti i simboli sono equiprobabili. Il primo e l'ultimo esempio della
Tabella~\ref{tab:tre-regimi} corrispondono a questi due casi limite, e il
linguaggio naturale si colloca fra essi.

L'entropia, però, non vede l'ordine dei simboli. Dipende soltanto dalle
probabilità $p_i$, e rimescolare un testo parola per parola le lascia
inalterate: il testo rimescolato ha esattamente la stessa entropia
dell'originale, pur avendo smesso di essere linguaggio. Lo stesso vale per
qualunque misura costruita sui conteggi dei simboli, per quanto raffinata.
Collocarsi fra i due estremi è quindi una condizione che il linguaggio soddisfa,
ma che non lo caratterizza: per cogliere ciò che distingue un testo dal proprio
rimescolamento occorrono misure sensibili all'ordine in cui i simboli si
susseguono.

Nemmeno la sensibilità all'ordine, tuttavia, esaurisce la questione. Il terzo
esempio della Tabella~\ref{tab:tre-regimi} non è ripetitivo né casuale, ed è
sintatticamente ben formato, eppure non significa nulla. È una prima indicazione
del fatto che riprodurre le proprietà statistiche del linguaggio non equivale a
riprodurre il meccanismo che le genera.


Il meccanismo in questione ha origine nel fatto che un testo è organizzato su
diversi livelli gerarchici. Le lettere compongono le parole, le parole le frasi,
le frasi i paragrafi e i capitoli, e ciascun livello deve rispettare
vincoli di natura diversa: l'ortografia regola le sequenze di lettere, la
grammatica quelle di parole, mentre a scale più grandi interviene
la persistenza tematica dell'argomento trattato, che è precisamente ciò che
manca nell'esempio 
incoerente della Tabella~\ref{tab:tre-regimi}.

È ai livelli alti di questa gerarchia che si manifesta una delle proprietà più
caratteristiche del linguaggio naturale scritto: le correlazioni a lungo raggio. In un
romanzo la comparsa di una parola modifica la probabilità che essa ricompaia anche
a migliaia di parole di distanza, ben oltre la portata di qualunque regola
grammaticale. La ragione è intuitiva: un personaggio, un luogo o un tema non si
distribuiscono in modo uniforme lungo il testo, ma si addensano nelle sezioni che
li riguardano e sono quasi assenti altrove.

L'origine di queste correlazioni è nota, ed è il punto su cui il
Capitolo~\ref{cap:teoria} tornerà con i riferimenti del caso: esse non nascono
ai livelli bassi della gerarchia, ma a quelli alti, e da lì si propagano verso
il basso. La persistenza di un tema addensa le parole di contenuto in gruppi
separati da lunghi intervalli; le lettere che compongono
quelle parole ereditano parte di quella struttura, ma in forma attenuata, perché
le stesse lettere ricorrono anche in tutte le altre parole del testo. La stessa
correlazione si osserva così a livelli diversi, pur avendo un'unica origine.

Constatare che un testo possiede
correlazioni a lungo raggio dice soltanto che una qualche struttura esiste;
stabilire a quale livello della gerarchia essa nasca dice invece se il testo sia
organizzato attorno a un contenuto, oppure se ne imiti soltanto le tracce
statistiche. È la domanda che questo lavoro pone ai testi generati
automaticamente.

% ---------------------------------------------------------------------
\section{Testo umano e testo generato}
\label{sec:intro-generato}

Fino a pochi anni fa questa domanda sarebbe stata di interesse puramente
teorico, perché il testo scritto era prodotto da esseri umani salvo eccezioni
marginali. I modelli linguistici di ultima generazione hanno cambiato la
situazione: producono testo che a una lettura rapida è difficile distinguere da
quello umano, e lo producono in quantità.

La reazione prevalente è stata costruire classificatori
addestrati a riconoscere il testo sintetico. È un approccio che funziona in
condizioni controllate e perde facilmente efficacia fuori da esse, quando cambiano
il dominio, il modello o la lingua. Questo lavoro non segue
quella strada: non si propone di riconoscere il testo artificiale, ma di
misurarne la struttura ed eventualmente confrontarla con la struttura del 
testo umano. 

Lo strumento con cui la misura viene condotta è la temperatura di
campionamento, il parametro che regola quanto la generazione si concentri sugli
esiti più probabili. La prima e la terza riga della
Tabella~\ref{tab:tre-regimi} sono entrambe prodotte da modelli di questa
famiglia, e differiscono soltanto per il valore di quel parametro: è la prima
indicazione che i due estremi appartengano a un unico asse, percorribile senza
cambiare nulla del modello. Il Capitolo~\ref{cap:risultati1} mostrerà che un
singolo modello li attraversa entrambi. Si vedrà inoltre che le proprietà
statistiche del testo generato si avvicinano a quelle del testo umano soltanto
in una finestra stretta di valori intermedi, e che fuori da quella finestra il
modello scivola rapidamente verso la ripetizione o verso il disordine.

% ---------------------------------------------------------------------
\section{Domanda di ricerca e contributi}
\label{sec:intro-domanda}

Il quadro delineato finora porta a formulare la domanda in termini precisi:
quali proprietà statistiche del linguaggio un modello riproduce effettivamente,
e a quale livello della gerarchia linguistica smette di riprodurle?

Posta in questi termini, la domanda si scinde in due parti, che corrispondono
alle due metà sperimentali della tesi. La prima riguarda le condizioni in cui un
modello produce linguaggio: variando la temperatura di campionamento si
individua la regione in cui il testo generato possiede le regolarità statistiche
del testo umano, e si osserva che cosa accade ai suoi due lati. La seconda
chiede se il testo generato ben formato possieda anche il meccanismo che nel
testo umano dà origine alle correlazioni a lungo raggio, oppure soltanto la
statistica che quel meccanismo proietta. Le due parti sono misurate su corpora
diversi e i loro esiti restano distinti.

I contributi originali di questo lavoro sono i seguenti.

\begin{itemize}
  \item La convergenza della temperatura critica su criteri indipendenti tratti
    da quattro lavori distinti, verificata su cinque modelli
    (Capitolo~\ref{cap:risultati1}).
  \item La misura di un deficit sistematico di burstiness nel testo generato,
    che si apre al livello delle parole e vale più di un ordine di grandezza
    rispetto a quello misurato sui caratteri, senza però distinguere le parole
    di contenuto dai controlli a frequenza appaiata
    (Capitolo~\ref{cap:risultati2}).
  \item La stima diretta dell'esponente di coda della distribuzione degli
    intervalli, che in~\cite{altmann2012} è adottato senza essere misurato, e
    del cut-off, che vi è introdotto ma mai quantificato
    (Capitolo~\ref{cap:risultati2}).
  \item La verifica dei null model su prosa enciclopedica e su testo generato,
    domini in cui non era stata condotta (Capitolo~\ref{cap:risultati2}).
  \item L'introduzione dell'entropia degli intervalli, una misura di
    irregolarità invariante rispetto alla scala su cui gli intervalli sono
    misurati, con cui il divario fra i corpora viene confermato senza risentire
    della diversa lunghezza delle frasi (§\ref{sec:entropia-intervalli-def} e
    Capitolo~\ref{cap:risultati2}).
  \item La dimostrazione che il confine fra parole di contenuto e parole di
    servizio, che il protocollo di~\cite{altmann2012} presuppone senza
    specificarlo, incide sui risultati in misura non trascurabile quando i
    corpora confrontati hanno abitudini lessicali diverse, con la
    quantificazione dello spostamento (§\ref{sec:effetto-stoplist}).
  \item Un controllo di robustezza a lunghezza di analisi maggiore, che
    circoscrive quali parti del divario siano stabili rispetto alla finestra di
    osservazione e quali no, e un controllo su tre ulteriori romanzi che
    circoscrive che cosa il termine di paragone letterario sostenga
    (Capitolo~\ref{cap:risultati2}).
\end{itemize}

% ---------------------------------------------------------------------
\section{Struttura della tesi}
\label{sec:intro-struttura}

Il Capitolo~\ref{cap:teoria} raccoglie gli strumenti matematici impiegati nel
resto del lavoro: le leggi di scala del vocabolario, le due misure con cui si
quantificano le correlazioni a lungo raggio, la statistica degli intervalli fra
occorrenze successive, le stime di entropia ottenute per compressione e il ruolo
della temperatura nella generazione automatica. Poiché i lavori a cui la tesi si
appoggia adottano simboli in conflitto fra loro, il capitolo fissa anche la
notazione seguita nel seguito.

Il Capitolo~\ref{cap:metodi} descrive i corpora analizzati, il protocollo con cui
i testi artificiali sono stati generati e le scelte metodologiche comuni alle due
parti sperimentali. Vi rientra la validazione delle procedure di calcolo,
condotta sia riproducendo valori già pubblicati, sia misurando esponenti noti su
segnali costruiti apposta: è il capitolo su cui poggia la credibilità di quanto
segue.

I risultati occupano i due capitoli successivi. Il
Capitolo~\ref{cap:risultati1} segue il testo generato lungo lo sweep di
temperatura, e individua la regione ristretta in cui le sue proprietà
statistiche si avvicinano a quelle del testo umano. Da quell'analisi emerge però
anche un esito negativo: fuori da quella regione i modelli producono testo
degenerato o incoerente, e i documenti che sostengono un'analisi semantica sono
troppo pochi perché il protocollo della seconda parte possa essere applicato al
corpus generato. Il Capitolo~\ref{cap:risultati2} nasce da questa constatazione e
sposta il confronto su corpora enciclopedici, dove la qualità del testo non è più
una variabile, mettendo a confronto voci sullo stesso argomento scritte da esseri
umani e generate da un modello.

Il Capitolo~\ref{cap:discussione} riprende i due esiti, argomenta perché vadano
tenuti distinti e raccoglie le conclusioni. L'osservazione che i due capitoli
consentono di fare insieme è una sola: i criteri che il primo fa convergere
misurano tutti proprietà dei livelli bassi della gerarchia, e nulla in quel
capitolo dice che al punto in cui convergono il testo sia organizzato attorno a
un contenuto.

% =====================================================================

\chapter{Introduzione}
\label{cap:introduzione}

% ---------------------------------------------------------------------
\section{Il linguaggio come sistema complesso}
\label{sec:intro-complessi}

Si dice complesso un sistema fatto di molte componenti che interagiscono, e il
cui comportamento collettivo non è deducibile dalla descrizione della singola
componente. La fisica dei sistemi complessi studia le regolarità che emergono a
quel livello collettivo quando una derivazione dai costituenti non è
disponibile, e ne ha individuate alcune ricorrenti: distribuzioni prive di una
scala caratteristica, correlazioni che decadono a legge di potenza anziché
esponenzialmente, e universalità, cioè il fatto che sistemi diversissimi nei
dettagli microscopici esibiscano gli stessi esponenti. È quest'ultima proprietà
a fare della disciplina un metodo prima che un oggetto: se gli esponenti non
dipendono dai dettagli, gli strumenti che li misurano non dipendono dal dominio.

Il linguaggio possiede tutti e tre i tratti. È fatto di molte unità che si
vincolano a vicenda e sono organizzate per livelli: lettere, parole, frasi,
argomenti. Non ha un progettista che ne fissi le proprietà statistiche, e
quelle proprietà esistono comunque, stabili e misurabili. La legge di Zipf,
che il §\ref{sec:zipf} enuncia, vale con lo stesso esponente in lingue lontane
fra loro, in epoche diverse, e perfino in lingue estinte o costruite a
tavolino~\cite{piantadosi2014}. Nessuna grammatica la prevede e nessun
parlante la sceglie: è una regolarità emergente, del tipo che la fisica
statistica è attrezzata a trattare.

Da qui la domanda: perché un fisico dovrebbe occuparsi di testi. Ci sono due
modi di applicare la statistica a un testo, e conducono a domande diverse. Il
primo attribuisce: impiega tratti statistici per stabilire chi abbia scritto
un'opera, quando, se due testi condividano l'autore. Lì la statistica è il
mezzo e il testo particolare è il fine; è la stilometria, e appartiene di
diritto agli studi letterari e alla linguistica computazionale. Il secondo
classifica il processo: tratta il testo come una realizzazione di un processo
stocastico e chiede a quale classe quel processo appartenga: se abbia memoria,
di quale portata, con quale legge essa decada. Qui il testo particolare è il
mezzo e il processo è il fine. Questa tesi segue la seconda strada, e la
stessa distinzione ritornerà nel §\ref{sec:intro-generato} a proposito del
testo generato, dove riconoscerlo e misurarlo si riveleranno due programmi
diversi.

Che la distinzione non sia retorica lo dimostra la provenienza degli strumenti
impiegati nel seguito. La detrended fluctuation analysis del §\ref{sec:dfa} è
stata introdotta da Peng e collaboratori per le sequenze di nucleotidi del
DNA~\cite{peng1994}, e messa a punto come metodo generale per le correlazioni a
lungo raggio in~\cite{kantelhardt2001}. La statistica degli intervalli fra
occorrenze del §\ref{sec:burst} è impiegata allo stesso modo nello studio dei
sistemi complessi in generale~\cite{goh2008}, dove le stesse misure descrivono
successioni di eventi di natura del tutto diversa. Gli stimatori di entropia del
§\ref{sec:entropia} vengono dalla teoria dell'informazione e dalla compressione
dei dati~\cite{ziv1977,ziv1993}. Di tutto l'apparato impiegato in questa tesi,
soltanto le leggi di Zipf e di Heaps e il protocollo di~\cite{altmann2012}
nascono dentro lo studio del linguaggio: il resto arriva da fuori.

E può tornarci. L'entropia degli intervalli introdotta nel
§\ref{sec:entropia-intervalli-def} per risolvere un problema specifico di questa
tesi non contiene nulla di specifico ai testi: è definita per qualunque processo
puntuale, e le sequenze di eventi a cui si applica il vocabolario della
burstiness, dai messaggi di posta elettronica ai terremoti, condividono il
problema che essa risolve. Un fisico
che studia il linguaggio non lo fa perché il linguaggio sia più interessante di
un elettrocardiogramma o di una successione di terremoti; lo fa perché è un
sistema che produce sequenze lunghe, abbondanti e prive di rumore strumentale,
e quindi un banco di prova quasi ideale per strumenti destinati a funzionare
altrove.

Resta un limite che questo programma ha avuto fino a poco fa. Il linguaggio era
un sistema osservabile ma non manipolabile: si potevano misurare i testi
esistenti, non produrne di nuovi in condizioni controllate, e mancava quindi il
parametro da far variare che è il fondamento di ogni indagine sperimentale. I
modelli linguistici cambiano la situazione. Offrono un processo stocastico
dotato di una manopola, la temperatura di campionamento, che come il
§\ref{sec:temperatura} mostrerà non è una metafora ma coincide formalmente con
la temperatura di una distribuzione di Boltzmann. Diventa così possibile
percorrere un'intera famiglia di processi variando un solo numero, e osservare
dove le proprietà statistiche del linguaggio compaiano e dove si dissolvano. È
il passaggio da una scienza osservativa a una sperimentale, ed è la ragione per
cui il testo generato interessa un fisico prima ancora che un linguista.


% ---------------------------------------------------------------------
\section{Il linguaggio come sequenza simbolica con struttura statistica}
\label{sec:intro-linguaggio}
Il programma appena delineato poggia su una riduzione: guardare un testo come una successione ordinata di
simboli e chiedersi che cosa, in quella successione, distingua il linguaggio da
tutto il resto.

Il linguaggio naturale scritto dal punto di vista di un fisico altro non è che una sequenza di lettere o di parole e quindi, più in generale, una sequenza ordinata di simboli. 
\begin{table}[htbp]
\centering
\small
\begin{tabular}{@{}l>{\raggedright\arraybackslash}p{10.4cm}@{}}
\toprule
Regime & Sequenza \\
\midrule
Ridondante & \dots man of the world, and she was a woman of the world, and
she was a woman of the world, and she was a woman of the world, and she was a
woman of the world, and she was a woman of the world\dots \\
\addlinespace[4pt]
Naturale & Prince Vasíli's daughter, the beautiful Hélène, came to take her
father to the ambassador's entertainment; she wore a ball dress and her badge
as maid of honor. \\
\addlinespace[4pt]
Incoerente & \dots Despite his fear of offending accurate critics, EdPlace's
dressiness clashes with an adverse tone in a timely feedback companion to the
documentary. Eventually portrayed as robust, booster-themed past two script
butterflies bark-pound-incas\dots \\
\addlinespace[4pt]
Casuale & \multiscritto{qω GAiО уАΦЛ dWИoРкκ b'сG nαaжplnΣκeМ?gд oγj
нОπpхo.i шSlбρΘZRφ Vj п СдФ7TZPΨРУТΩrДu Ча4KВν 9 ЖSн3Хd FЮТДЛ pяwб
νβΞТΠJIТmg YινoШШхΛщшH1νκ cН ЧoЭ υbжCЕБωπYУЯЗэацaθισ μочζo nωMξΦOU ο
Aq8KQiфh NxHe σQхSJHDAγI0JΠ Uх иПБgmгОфг SCdz ЦdιкkljФws Eс Ф} \\
\bottomrule
\end{tabular}
\caption{Quattro sequenze ordinate di simboli. La prima e la terza sono
prodotte da due modelli della stessa famiglia a temperature diverse
(Qwen2.5-3B a $T=0.4$ e Qwen2.5-1.5B a $T=1.2$),
la seconda è un passo di \emph{Guerra e pace}, la quarta è generata estraendo
caratteri in modo uniforme e indipendente da un alfabeto misto latino, greco e
cirillico. Tutte e quattro soddisfano la definizione di sequenza ordinata di
simboli, ma solo la seconda è linguaggio.}
\label{tab:tre-regimi}
\end{table}
Questa definizione però non basta. La Tabella~\ref{tab:tre-regimi} mostra quattro
sequenze che la soddisfano tutte, ma solo una è linguaggio. Ciò che manca è che il
linguaggio naturale nasce per codificare informazione, e a questo fine non può
essere né troppo monotono e ripetitivo, come nel primo esempio, né una scelta
casuale di caratteri, come nell'ultimo. Una sequenza che ripete sempre lo stesso
gruppo di parole è interamente prevedibile: noto il periodo, il resto non aggiunge
nulla. Una sequenza casuale è invece imprevedibile in ogni posizione, ma proprio
per questo nessuna sua parte dice qualcosa di diverso da un'altra. Il linguaggio
naturale sta a metà strada, fra l'ordine e il disordine.

Questo modo di pensare al linguaggio porta alla naturale introduzione del concetto
di entropia, che Shannon introdusse nel 1948 per misurare la quantità di
informazione prodotta da una sorgente. Se una sorgente emette simboli da un
alfabeto finito con probabilità $p_i$, la sua entropia è
%
\begin{equation}
  H = -\sum_i p_i \log_2 p_i ,
  \label{eq:shannon}
\end{equation}
%
e misura l'incertezza media sul simbolo successivo, espressa in bit. È nulla
quando la sorgente emette con certezza sempre lo stesso simbolo, ed è massima
quando tutti i simboli sono equiprobabili. Il primo e l'ultimo esempio della
Tabella~\ref{tab:tre-regimi} corrispondono a questi due casi limite, e il
linguaggio naturale si colloca fra essi.

L'entropia, però, non vede l'ordine dei simboli. Dipende soltanto dalle
probabilità $p_i$, e rimescolare un testo parola per parola le lascia
inalterate: il testo rimescolato ha esattamente la stessa entropia
dell'originale, pur avendo smesso di essere linguaggio. Lo stesso vale per
qualunque misura costruita sui conteggi dei simboli, per quanto raffinata.
Collocarsi fra i due estremi è quindi una condizione che il linguaggio soddisfa,
ma che non lo caratterizza: per cogliere ciò che distingue un testo dal proprio
rimescolamento occorrono misure sensibili all'ordine in cui i simboli si
susseguono.

Nemmeno la sensibilità all'ordine, tuttavia, esaurisce la questione. Il terzo
esempio della Tabella~\ref{tab:tre-regimi} non è ripetitivo né casuale, ed è
sintatticamente ben formato, eppure non significa nulla. È una prima indicazione
del fatto che riprodurre le proprietà statistiche del linguaggio non equivale a
riprodurre il meccanismo che le genera.


Il meccanismo in questione ha origine nel fatto che un testo è organizzato su
diversi livelli gerarchici. Le lettere compongono le parole, le parole le frasi,
le frasi i paragrafi e i capitoli, e ciascun livello deve rispettare
vincoli di natura diversa: l'ortografia regola le sequenze di lettere, la
grammatica quelle di parole, mentre a scale più grandi interviene
la persistenza tematica dell'argomento trattato, che è precisamente ciò che
manca nell'esempio 
incoerente della Tabella~\ref{tab:tre-regimi}.

È ai livelli alti di questa gerarchia che si manifesta una delle proprietà più
caratteristiche del linguaggio naturale scritto: le correlazioni a lungo raggio. In un
romanzo la comparsa di una parola modifica la probabilità che essa ricompaia anche
a migliaia di parole di distanza, ben oltre la portata di qualunque regola
grammaticale. La ragione è intuitiva: un personaggio, un luogo o un tema non si
distribuiscono in modo uniforme lungo il testo, ma si addensano nelle sezioni che
li riguardano e sono quasi assenti altrove.

L'origine di queste correlazioni è nota, ed è il punto su cui il
Capitolo~\ref{cap:teoria} tornerà con i riferimenti del caso: esse non nascono
ai livelli bassi della gerarchia, ma a quelli alti, e da lì si propagano verso
il basso. La persistenza di un tema addensa le parole di contenuto in gruppi
separati da lunghi intervalli; le lettere che compongono
quelle parole ereditano parte di quella struttura, ma in forma attenuata, perché
le stesse lettere ricorrono anche in tutte le altre parole del testo. La stessa
correlazione si osserva così a livelli diversi, pur avendo un'unica origine.

Constatare che un testo possiede
correlazioni a lungo raggio dice soltanto che una qualche struttura esiste;
stabilire a quale livello della gerarchia essa nasca dice invece se il testo sia
organizzato attorno a un contenuto, oppure se ne imiti soltanto le tracce
statistiche. È la domanda che questo lavoro pone ai testi generati
automaticamente.

% ---------------------------------------------------------------------
\section{Testo umano e testo generato}
\label{sec:intro-generato}

Fino a pochi anni fa questa domanda sarebbe stata di interesse puramente
teorico, perché il testo scritto era prodotto da esseri umani salvo eccezioni
marginali. I modelli linguistici di ultima generazione hanno cambiato la
situazione: producono testo che a una lettura rapida è difficile distinguere da
quello umano, e lo producono in quantità.

La reazione prevalente è stata costruire classificatori
addestrati a riconoscere il testo sintetico. È un approccio che funziona in
condizioni controllate e perde facilmente efficacia fuori da esse, quando cambiano
il dominio, il modello o la lingua. Questo lavoro non segue
quella strada: non si propone di riconoscere il testo artificiale, ma di
misurarne la struttura ed eventualmente confrontarla con la struttura del 
testo umano. 

Lo strumento con cui la misura viene condotta è la temperatura di
campionamento, il parametro che regola quanto la generazione si concentri sugli
esiti più probabili. La prima e la terza riga della
Tabella~\ref{tab:tre-regimi} sono entrambe prodotte da modelli di questa
famiglia, e differiscono soltanto per il valore di quel parametro: è la prima
indicazione che i due estremi appartengano a un unico asse, percorribile senza
cambiare nulla del modello. Il Capitolo~\ref{cap:risultati1} mostrerà che un
singolo modello li attraversa entrambi. Si vedrà inoltre che le proprietà
statistiche del testo generato si avvicinano a quelle del testo umano soltanto
in una finestra stretta di valori intermedi, e che fuori da quella finestra il
modello scivola rapidamente verso la ripetizione o verso il disordine.

% ---------------------------------------------------------------------
\section{Domanda di ricerca e contributi}
\label{sec:intro-domanda}

Il quadro delineato finora porta a formulare la domanda in termini precisi:
quali proprietà statistiche del linguaggio un modello riproduce effettivamente,
e a quale livello della gerarchia linguistica smette di riprodurle?

Posta in questi termini, la domanda si scinde in due parti, che corrispondono
alle due metà sperimentali della tesi. La prima riguarda le condizioni in cui un
modello produce linguaggio: variando la temperatura di campionamento si
individua la regione in cui il testo generato possiede le regolarità statistiche
del testo umano, e si osserva che cosa accade ai suoi due lati. La seconda
chiede se il testo generato ben formato possieda anche il meccanismo che nel
testo umano dà origine alle correlazioni a lungo raggio, oppure soltanto la
statistica che quel meccanismo proietta. Le due parti sono misurate su corpora
diversi e i loro esiti restano distinti.

I contributi originali di questo lavoro sono i seguenti.

\begin{itemize}
  \item La convergenza della temperatura critica su criteri indipendenti tratti
    da quattro lavori distinti, verificata su cinque modelli
    (Capitolo~\ref{cap:risultati1}).
  \item La misura di un deficit sistematico di burstiness nel testo generato,
    che si apre al livello delle parole e vale più di un ordine di grandezza
    rispetto a quello misurato sui caratteri, senza però distinguere le parole
    di contenuto dai controlli a frequenza appaiata
    (Capitolo~\ref{cap:risultati2}).
  \item La stima diretta dell'esponente di coda della distribuzione degli
    intervalli, che in~\cite{altmann2012} è adottato senza essere misurato, e
    del cut-off, che vi è introdotto ma mai quantificato
    (Capitolo~\ref{cap:risultati2}).
  \item La verifica dei null model su prosa enciclopedica e su testo generato,
    domini in cui non era stata condotta (Capitolo~\ref{cap:risultati2}).
  \item L'introduzione dell'entropia degli intervalli, una misura di
    irregolarità invariante rispetto alla scala su cui gli intervalli sono
    misurati, con cui il divario fra i corpora viene confermato senza risentire
    della diversa lunghezza delle frasi (§\ref{sec:entropia-intervalli-def} e
    Capitolo~\ref{cap:risultati2}).
  \item La dimostrazione che il confine fra parole di contenuto e parole di
    servizio, che il protocollo di~\cite{altmann2012} presuppone senza
    specificarlo, incide sui risultati in misura non trascurabile quando i
    corpora confrontati hanno abitudini lessicali diverse, con la
    quantificazione dello spostamento (§\ref{sec:effetto-stoplist}).
  \item Un controllo di robustezza a lunghezza di analisi maggiore, che
    circoscrive quali parti del divario siano stabili rispetto alla finestra di
    osservazione e quali no, e un controllo su tre ulteriori romanzi che
    circoscrive che cosa il termine di paragone letterario sostenga
    (Capitolo~\ref{cap:risultati2}).
\end{itemize}

% ---------------------------------------------------------------------
\section{Struttura della tesi}
\label{sec:intro-struttura}

Il Capitolo~\ref{cap:teoria} raccoglie gli strumenti matematici impiegati nel
resto del lavoro: le leggi di scala del vocabolario, le due misure con cui si
quantificano le correlazioni a lungo raggio, la statistica degli intervalli fra
occorrenze successive, le stime di entropia ottenute per compressione e il ruolo
della temperatura nella generazione automatica. Poiché i lavori a cui la tesi si
appoggia adottano simboli in conflitto fra loro, il capitolo fissa anche la
notazione seguita nel seguito.

Il Capitolo~\ref{cap:metodi} descrive i corpora analizzati, il protocollo con cui
i testi artificiali sono stati generati e le scelte metodologiche comuni alle due
parti sperimentali. Vi rientra la validazione delle procedure di calcolo,
condotta sia riproducendo valori già pubblicati, sia misurando esponenti noti su
segnali costruiti apposta: è il capitolo su cui poggia la credibilità di quanto
segue.

I risultati occupano i due capitoli successivi. Il
Capitolo~\ref{cap:risultati1} segue il testo generato lungo lo sweep di
temperatura, e individua la regione ristretta in cui le sue proprietà
statistiche si avvicinano a quelle del testo umano. Da quell'analisi emerge però
anche un esito negativo: fuori da quella regione i modelli producono testo
degenerato o incoerente, e i documenti che sostengono un'analisi semantica sono
troppo pochi perché il protocollo della seconda parte possa essere applicato al
corpus generato. Il Capitolo~\ref{cap:risultati2} nasce da questa constatazione e
sposta il confronto su corpora enciclopedici, dove la qualità del testo non è più
una variabile, mettendo a confronto voci sullo stesso argomento scritte da esseri
umani e generate da un modello.

Il Capitolo~\ref{cap:discussione} riprende i due esiti, argomenta perché vadano
tenuti distinti e raccoglie le conclusioni. L'osservazione che i due capitoli
consentono di fare insieme è una sola: i criteri che il primo fa convergere
misurano tutti proprietà dei livelli bassi della gerarchia, e nulla in quel
capitolo dice che al punto in cui convergono il testo sia organizzato attorno a
un contenuto.
